\chapter{Diseño metodológico}
\label{cap:metodologia}

Cómo se baja OE1--OE4 a pasos. Enfoque, fuentes, técnicas y dónde se puede romper la validez.

\section{Tipo y enfoque metodológico}
Es investigación tecnológica (\textit{design science}): se diseña un artefacto ---el modelo
S/E/C $\rightarrow$ \AlignmentScore{}--- que convierte YAML en números. No hay sujetos humanos.
Las entradas son contratos OpenAPI y extractos BIAN (§1.9).

Guía~1 deja el diseño. El software y el experimento a escala son Tesis~2 (§1.6). Aquí se dice de
dónde salen los datos, qué par se mide y qué factores se podrían contrastar.


\section{Procedimiento general}
Cinco fases, las del Anexo~C. Primero se identifica y versiona el YAML y el extracto BIAN
(§4.3); después se parsea. Aquí se escribe el diseño. Correrlo a escala es §1.6.

\begin{enumerate}
  \item \textbf{Fase 1 --- Diseñar el modelo:} definir entradas, scores S/E/C, fórmula de
        agregación y umbrales de clasificación (Cap.~1 §1.4; Anexo~C).
  \item \textbf{Fase 2 --- Adquirir y normalizar artefactos (OE1):} obtener contratos OpenAPI y
        extractos BIAN (§4.3); parsear y producir representaciones comparables (tabular, grafo o RDF
        según factor elegido).
  \item \textbf{Fase 3 --- Emparejar estructuras (OE2):} calcular correspondencias y
        \StructuralScore{} (E).
  \item \textbf{Fase 4 --- Calcular similitud semántica (OE3):} obtener \SemanticScore{} (S).
  \item \textbf{Fase 5 --- Integrar y clasificar (OE4):} calcular \CoverageScore{} (C),
        \AlignmentScore{} y clasificación por tipologías (Tabla~\ref{tab:tipologias}).
\end{enumerate}

\subsection{Relación objetivos, fases y procedimiento}
\label{sec:objetivos-metodologia}

La \Cref{tab:oe-metodologia} pone cada OE en una fase, con un verbo y un producto.

\begin{table}[H]
  \centering
  \caption{Objetivos específicos vinculados al procedimiento metodológico}
  \label{tab:oe-metodologia}
  \small
  \begin{tabular}{@{}p{0.06\textwidth}p{0.14\textwidth}p{0.22\textwidth}p{0.24\textwidth}p{0.24\textwidth}@{}}
    \toprule
    \textbf{OE} & \textbf{Fase} & \textbf{Acción (verbo)} & \textbf{Entrada} & \textbf{Producto} \\
    \midrule
    --- &
    Fase 1 &
    \textbf{Diseñar} entradas, scores y umbrales &
    Requisitos del §1.4; Anexo C &
    Modelo S/E/C documentado \\
    \textbf{OE1} &
    Fase 2 &
    \textbf{Adquirir} y \textbf{normalizar} contrato y SD &
    OpenAPI + extracto BIAN (§4.3) &
    Representaciones comparables \\
    \textbf{OE2} &
    Fase 3 &
    \textbf{Emparejar} estructuras y calcular E &
    Representaciones OE1 &
    Matriz de correspondencias + E \\
    \textbf{OE3} &
    Fase 4 &
    \textbf{Calcular} similitud semántica S &
    Representaciones OE1 &
    Matriz semántica + S \\
    \textbf{OE4} &
    Fase 5 &
    \textbf{Integrar} S, E, C y clasificar &
    E, S, representaciones OE1 &
    C, \AlignmentScore{}, tipologías \\
    \bottomrule
  \end{tabular}
\end{table}

Las técnicas por OE están en la \Cref{sec:tecnicas-herramientas}. Las filas, en el Anexo~C.


La \Cref{fig:proc-estructural} ilustra la correspondencia jerárquica para E; la cadena completa
aparece en la \Cref{fig:cadena-objetivos} (Cap.~2).

\setcounter{figure}{6}
\begin{figure}[H]
  \centering
  \figCapSeven
  \caption{Correspondencia jerárquica para el cálculo de \StructuralScore{} (E). Fuente: elaboración
  propia; ejemplo ilustrativo.}
  \label{fig:proc-estructural}
\end{figure}

\subsection{Protocolo de validación empírica (fuera de alcance)}

Cuando se corra, hará falta un repositorio de contratos, el parser, un registro de instancia y
versión, un criterio de muestra (no un YAML suelto) y ---si se compara contra un humano--- un
criterio de par ``correcto'' escrito, no improvisado. Ley 29733: sin PII. Precisión contra esa
referencia y falsos pares se reportan ahí, no en Guía~1.


\section{Fuentes de datos}
\subsection{Adquisición de datos}

No hay encuesta. Hay YAML/JSON OpenAPI y material BIAN. Cada instancia lleva un id, p.\ ej.\
\texttt{\{fuente\}\_\{service-domain\}\_\{version\}}.

\begin{table}[H]
  \centering
  \caption{Adquisición de datos --- diseño metodológico}
  \label{tab:adquisicion-datos}
  \small
  \begin{tabular}{@{}p{0.26\textwidth}p{0.66\textwidth}@{}}
    \toprule
    \textbf{Elemento} & \textbf{En este estudio} \\
    \midrule
    Sistema / entorno &
    Sector bancario; arquitecturas API + modelos de referencia BIAN \\
    Objeto de estudio &
    Contrato OpenAPI bancario completo (Cap.~1 §1.4) \\
    Unidad de análisis &
    Par comparable endpoint/\allowbreak schema $\leftrightarrow$ \Behavior{}/\BusinessObject{} \\
    Señal / dato &
    Archivos OpenAPI (YAML/JSON) + extracto del \ServiceDomain{} BIAN \\
    Instrumento de medición &
    Pipeline del modelo (normalización + cálculo S, E, C); \textit{parsing} OpenAPI \\
    Parámetros controlados &
    Versión OpenAPI~3.x; dominio SD; idioma de descripciones; herramienta de \textit{parsing} \\
    Parámetros no controlados &
    Calidad redaccional del contrato; madurez BIAN del emisor; completitud del SD publicado \\
    Factores (Cap.~2) &
    Esquema intermedio; estrategia de matching; técnica semántica; pesos $\alpha,\beta,\gamma$ \\
    Ratios (resultados) &
    \StructuralScore{}, \SemanticScore{}, \CoverageScore{}, \AlignmentScore{}, clasificación \\
    Formato almacenamiento &
    YAML/JSON versionado en repositorio; metadatos CSV/JSON por instancia \\
    Tamaño de muestra &
    Por definir en validación empírica posterior; ilustración actual: Anexo~F (Payment Initiation) \\
    \bottomrule
  \end{tabular}
\end{table}


\subsection{Criterios de inclusión (muestra empírica)}

Para una validación empírica posterior se considerarán contratos OpenAPI~3.x del sector bancario
emparejados con un \ServiceDomain{} BIAN declarado en la documentación del caso, con:

\begin{itemize}
  \item especificación parseable sin errores de sintaxis;
  \item al menos una operación y un esquema en \texttt{components/schemas};
  \item referencia BIAN disponible en formato legible (catálogo, extracto o documentación oficial);
  \item ausencia de datos personales identificables (Ley N°~29733): preferir ejemplos públicos,
        sintéticos o anonimizados (Anexo~F).
\end{itemize}

La población es todo par OpenAPI $\leftrightarrow$ SD que cumpla eso. La muestra, un recorte
(pagos, cuentas, préstamos\ldots). Guía~1 no clava un piloto único.


\section{Técnicas y herramientas}
\label{sec:tecnicas-herramientas}
Por objetivo, tomadas del §\ref{sec:estado-arte} y de los factores del Capítulo~\ref{cap:objetivos}.
No hay aprendizaje supervisado: S y E salen de reglas de \textit{schema matching} y de similitud,
no de un modelo etiquetado.

\begin{itemize}
  \item \textbf{OE1:} \textit{parsing} OpenAPI~3.x; extracción de \textit{paths}, operaciones y
        esquemas; representación intermedia alineada con BIAN (Mainas et al., 2023 como referencia
        de enriquecimiento semántico).
  \item \textbf{OE2:} \textit{schema matching} según taxonomía Shvaiko y Euzenat (2005): nombre,
        ruta, profundidad de esquema, operación HTTP.
  \item \textbf{OE3:} similitud semántica por embeddings y/o ontología (El Bayadh et al., 2015;
        Chandrasekaran y Mago, 2022).
  \item \textbf{OE4:} agregación ponderada $\alpha S + \beta E + \gamma C$ y clasificación por
        umbrales fijos (§1.4).
\end{itemize}

Herramientas previstas: Python para prototipo, validadores OpenAPI, bibliotecas de similitud textual.
La implementación concreta queda fuera del alcance del presente plan (§\ref{sec:alcances}).


\section{Métricas de evaluación}
Las definiciones operativas están en el Cap.~1 §1.4 y en el Anexo~D: \StructuralScore{} y
\SemanticScore{} en $[0,1]$, \CoverageScore{} como proporción de elementos BIAN cubiertos, y
\AlignmentScore{} con los cortes Alta ($\geq 0{,}80$), Media, Baja y Nula.

Los indicadores por fase están en la matriz del Anexo~C y en la \Cref{tab:anexo-d}.


\section{Casos de estudio}
Los casos se eligen por dominio, no porque el Anexo~F sea ``el piloto''. Ese anexo
(\textit{Payment Initiation}) solo muestra el formato: un \texttt{POST /payment-orders} frente a
un \Behavior{} \texttt{Initiate}, y así.

Cada caso debería dejar: id de instancia, YAML, extracto BIAN, tabla de pares y los cuatro
números (S, E, C, \AlignmentScore{}). Cuántos casos y cuáles se corren no se decide en Guía~1
(§1.6).


\section{Criterios de validez}

Cuatro anclas, ninguna estadística sobre personas: el Anexo~C (problema $\rightarrow$ producto),
que otro pueda repetir los pasos, que Cap.~1 / Cap.~3 / Anexo~D no se contradigan, y el recorte
de §1.6.

\section{Amenazas a la validez}

Sin software en esta Guía, la amenaza interna es un parser que se coma mal un \texttt{allOf} o
un S sesgado porque \texttt{Retrieve} y ``consultar'' no se normalizaron en OE1.

Hacia afuera: 0,80 / 0,60 / 0,40 son convención de diseño. Un SD de crédito puede no
comportarse como uno de pagos; otra \textit{release} BIAN tampoco.

El constructo flaquea si un \textit{schema} crítico queda mal emparejado y C lo diluye: el
\AlignmentScore{} sale ``Medio'' y negocio sigue roto.

Nada de esto se generaliza al sector con un YAML. Hace falta el criterio de inclusión de §4.3 y
reportar la instancia junto al score. Eso es Tesis~2.
